\documentclass[10pt]{article}
\usepackage[a4paper,margin=2cm]{geometry}
\usepackage{amsmath,amssymb,amsthm}
\usepackage{listings}
\usepackage{xcolor}
\usepackage{hyperref}
\lstset{
  basicstyle=\ttfamily\small,
  breaklines=true,
  showstringspaces=false,
  columns=fullflexible,
}

\title{Plan: \textbf{HymeKo-Nagare} (HymeKo-Flow) ---\\
  A Paradigm-Native Dataflow ML Framework for\\
  Signed-Hypergraph Learning}
\author{}
\date{2026-05-11}

\begin{document}
\maketitle

\section*{Etymology}

\textbf{HymeKo-Nagare} (流れ): \emph{nagare} is Japanese for ``flow / current
/ stream.'' Combined with the HymeKo brand and pronounced cleanly in
Hungarian (\emph{na-ga-ré}). Public-facing English alias:
\textbf{HymeKo-Flow}.

\section*{Scope and motivation}

PyTorch, JAX, Burn, Candle, and dfdx all share a single computational
paradigm: \emph{dense N-d tensors as the universal data type, composed
into a directed-acyclic computational graph, differentiated by
reverse-mode autograd over that graph.}

This paradigm is well-suited to convolutional / transformer / dense
models where computation is naturally expressed as a sequence of
tensor operations with a batch dimension that is the only embarrassingly
parallel axis within a layer.

\textbf{It is NOT a good fit for signed-hypergraph learning.} Our
forward computation factors as
\[
  F(G, X) = \sum_{c \in \mathcal{C}} f_\theta(c, X),
  \qquad
  \frac{\partial F}{\partial \theta} = \sum_{c \in \mathcal{C}}
      \frac{\partial f_\theta(c, X)}{\partial \theta},
\]
where $\mathcal{C}$ is the cycle pool (or walk pool, or signed-hyperedge
pool) and $f_\theta$ is a small per-cycle aggregator with shared
parameters $\theta$. Both forward and gradient are \emph{sums of
independent terms} over $|\mathcal{C}|$. The per-cycle computation is
embarrassingly parallel; the gradient accumulation is a commutative
reduction.

This is the structure of a \textbf{MapReduce / dataflow} computation,
not of a tensor-op DAG. Expressing it as the latter forces:
\begin{itemize}
  \item materialising cycle pools as $(|\mathcal{C}|, k, d)$ dense
    tensors (memory waste at $k \ll d$);
  \item per-op autograd-graph nodes that prevent operator fusion
    across cycles;
  \item serialisation through PyTorch's eager-mode Python dispatcher
    (the source of the $22\times$ Rust-SIMD speedup we measured
    today over PyTorch eager FIR on Bitcoin OTC).
\end{itemize}

HymeKo-Nagare is the framework that takes signed-hypergraph learning's
natural decomposition seriously: \emph{operators are
$(\text{forward}, \text{backward})$ pairs over SoA cycle pools,
composed by direct function call, parallelised by cycle, executed on
CPU via Rayon or GPU via wgpu/cudarc, with no autograd graph.}

\section*{Affected files}

\textbf{New top-level crate}: \texttt{hymeko\_nagare/}
\begin{lstlisting}
hymeko_nagare/
├── Cargo.toml
├── src/
│   ├── lib.rs                   pub crate exports
│   ├── types.rs                 SoA CyclePool, Params, Grads
│   ├── ops/
│   │   ├── mod.rs               public Op trait + registry
│   │   ├── fir.rs               signed-cycle FIR fwd+bwd
│   │   ├── spline.rs            Catmull-Rom signed-cycle fwd+bwd
│   │   ├── scatter.rs           scatter-mean + scatter-sum fwd+bwd
│   │   ├── mlp.rs               edge predictor fwd+bwd
│   │   ├── loss.rs              BCE-with-logits fwd+bwd
│   │   └── adam.rs              optimizer step
│   ├── backend/
│   │   ├── mod.rs               Backend trait
│   │   ├── cpu_rayon.rs         Rayon parallel-by-cycle
│   │   └── gpu_wgpu.rs          wgpu compute-shader path
│   ├── grad_accum.rs            atomic-add gradient accumulator
│   └── training.rs              forward+backward+step composition
└── tests/                       per-op fwd+bwd unit tests
\end{lstlisting}

\textbf{Touched (additive only)}: \texttt{hymeko\_py/src/nagare.rs}
(PyO3 bindings so PyTorch-using callers can drop the Rust path in
during inference). PyTorch itself is untouched.

\section*{CORE.YAML items touched}

None. New crate, additive throughout. Existing flat HSiKAN pipeline
in \texttt{signedkan\_wip/src/run\_final\_cell.py} stays as the
training/research path. HymeKo-Nagare initially targets \emph{inference}
deployment only; training migration is a Stage 2 follow-up (see
``Sequencing'' below).

\section*{Mathematical foundations}

\subsection*{Cycle-additive operator algebra}

Define the cycle-additive operator class:
\[
  \mathcal{A} = \big\{ A : \mathcal{C} \times \mathbb{R}^{N \times d}
                       \to \mathbb{R}^p \mid
                       A(\mathcal{C}, X) = \sum_{c} a_\theta(c, X) \big\}.
\]
Properties:
\begin{enumerate}
  \item \textbf{Closure under sum and composition with cycle-additive
    aggregators.} If $A_1, A_2 \in \mathcal{A}$ and
    $g : \mathbb{R}^p \to \mathbb{R}^q$ is per-cycle, then
    $g \circ A \in \mathcal{A}$ and $A_1 + A_2 \in \mathcal{A}$.
  \item \textbf{Commutativity of evaluation order.} Output is
    invariant under any permutation of $\mathcal{C}$ since sum is
    commutative. (Up to float-summation order; mitigated with
    Kahan summation when needed, per CLAUDE.md \S5 numerical
    stability rules.)
  \item \textbf{Gradient additivity.} For any $\theta$,
    $\partial A / \partial \theta = \sum_c \partial a_\theta(c, X) / \partial \theta$.
    Gradient accumulation is a commutative reduction; lockless
    atomic-add on shared $\theta_{\text{grad}}$ is sound when the
    underlying float-add is atomic (\texttt{std::sync::atomic::AtomicU32}
    via bit-cast \texttt{f32} $\leftrightarrow$ \texttt{u32} with
    compare-exchange loop).
  \item \textbf{Scatter-mean is a cycle-additive reduce.} The
    per-vertex aggregation
    $H[v] = \sum_{c: v \in c} a_\theta(c, X) / |\{c: v \in c\}|$
    is a vertex-keyed reduction over the cycle-map outputs. Backward:
    each $a_\theta(c, X)$ receives gradient
    $\sum_{v \in c} \partial L / \partial H[v] \cdot
    |\{c': v \in c'\}|^{-1}$, also cycle-additive.
\end{enumerate}

\subsection*{Per-primitive forward + backward kernels}

For each primitive we need both kernels in pure Rust:

\paragraph{FIR forward}
$\mathrm{out}_c[j] = \sum_{i=0}^{k-1} k_{\sigma_i}[i] \cdot X[v_i][j]$
--- already shipped (see \texttt{hymeko\_graph::spine}).

\paragraph{FIR backward}
\[
  \frac{\partial L}{\partial k_{+1}[i]} = \sum_{c} \mathbb{1}[\sigma_i = +1]
    \cdot X[v_i] \cdot \frac{\partial L}{\partial \mathrm{out}_c},
  \quad
  \frac{\partial L}{\partial X[v]} = \sum_{c, i: v_i = v}
    k_{\sigma_i}[i] \cdot \frac{\partial L}{\partial \mathrm{out}_c}.
\]
Both Rayon-parallel over cycles; the $X$ gradient uses atomic add
because multiple cycles can write to the same vertex.

\paragraph{Catmull-Rom spline forward}
Per the existing PyTorch reference in \texttt{signedkan\_wip/src/splines.py};
port to pure Rust with the same numerical conventions, verified to
$\le 10^{-7}$ vs PyTorch.

\paragraph{Catmull-Rom spline backward}
Already have closed-form gradients (Triton kernel in
\texttt{signedkan\_wip/src/triton\_fused\_backward.py} hands them off
to the kernel); port the math to Rust.

\paragraph{Edge-predictor MLP forward + backward}
2-layer GELU MLP; standard. $\sim$200 LOC including GELU + chain rule.

\paragraph{BCE-with-logits forward + backward}
\[
  L = \tfrac{1}{N}\sum (\max(s, 0) - s \cdot y + \log(1 + e^{-|s|})),
  \quad
  \partial L / \partial s = \tfrac{1}{N}(\sigma(s) - y).
\]

\paragraph{Adam step}
\[
  m \leftarrow \beta_1 m + (1-\beta_1) g, \quad
  v \leftarrow \beta_2 v + (1-\beta_2) g^2, \quad
  \theta \leftarrow \theta - \alpha \hat m / (\sqrt{\hat v} + \epsilon).
\]
$\sim$50 LOC pure Rust per parameter group.

\section*{Backend abstraction}

\begin{lstlisting}
trait Backend {
    /// FIR forward over a cycle pool.
    fn fir_forward(
        cycles: &CyclePool,
        features: &[f32], d: usize,
        params: &FIRParams,
    ) -> Vec<f32>;
    fn fir_backward(
        cycles: &CyclePool,
        features: &[f32], d: usize,
        params: &FIRParams,
        grad_out: &[f32],
    ) -> (Vec<f32>, FIRGrads);  // (grad_features, grad_params)
    // ... same shape for spline, scatter, mlp, etc.
}

struct CpuRayonBackend;
struct WgpuBackend;       // GPU via wgpu compute shaders (portable)
struct CudarcBackend;     // GPU via cudarc (NVIDIA-only, faster)
\end{lstlisting}

\textbf{Same operator API across backends.} The execution scheduler
picks a backend per-op based on input size (small graph: CPU Rayon;
large graph: GPU). This is the \emph{total parallelism} story: every
backend exposes the same MapReduce primitive over cycles, and the
choice is just where the map lands.

\section*{Test strategy (CLAUDE.md \S3)}

For each $(\text{op}, \text{backend})$ pair:
\begin{itemize}
  \item \textbf{Forward parity} test against PyTorch reference to
    $\le 10^{-7}$ on random fixtures (5 random shapes/inputs).
  \item \textbf{Numerical gradient} test:
    $\partial L / \partial \theta_{ij}$ via central differences
    matches the analytical backward to $\le 10^{-3}$.
  \item \textbf{Atomic-add correctness}: 5 Rayon threads racing on
    the same $X$-gradient accumulator produce the same result as
    sequential evaluation (mod float-summation order, mitigated by
    Kahan).
  \item \textbf{Backend equivalence}: CPU and GPU produce the same
    output to $\le 10^{-5}$ on shared random fixture.
\end{itemize}

\textbf{Smoke}: train an FIR-only signed-cycle model on Bitcoin OTC
1 epoch in HymeKo-Nagare, verify loss decreases monotonically, AUC
$> 0.55$ (better than random), peak memory $< 4$ GB. Same smoke gate
as CPML had.

\section*{Performance budget}

\begin{itemize}
  \item Bitcoin OTC inference: $\le 1$ ms / iter at $d=32, k=3$
    (already achieved by spine alone).
  \item Bitcoin OTC training (FIR-only): $\le 5$ s / epoch on CPU.
  \item Slashdot training: $\le 60$ s / epoch on CPU; $\le 10$ s on
    GPU via wgpu / cudarc.
  \item Epinions training: $\le 300$ s / epoch on CPU; $\le 60$ s on
    GPU.
  \item Peak RAM: $\le 4 \cdot d \cdot |V| + 4 \cdot k \cdot |\mathcal{C}|$
    bytes (the SoA storage) plus $O(p)$ parameters + grads.
\end{itemize}

\section*{Sequencing (each stage is its own plan-then-implement gate)}

\begin{enumerate}
  \item \textbf{Stage 1 (inference only)}: FIR forward + scatter-mean +
    MLP forward + edge predictor forward, no backward, no optimizer.
    Drop-in inference replacement for trained HSiKAN models. CPU
    Rayon backend only. $\sim$1 month.
  \item \textbf{Stage 2 (training, CPU)}: add backward for every Stage-1
    op; Adam optimizer; training loop; atomic-grad-accumulator.
    Smoke-validate on Bitcoin OTC vs PyTorch reference. $\sim$2 months.
  \item \textbf{Stage 3 (GPU backend)}: wgpu compute shaders for FIR
    fwd+bwd, scatter, MLP. Cudarc as faster NVIDIA-specific path.
    Backend selection at op-call time. $\sim$2 months.
  \item \textbf{Stage 4 (full primitives)}: Catmull-Rom spline fwd+bwd
    in pure Rust (port from Triton); $\alpha_k$ multi-arity routing;
    attention ops if needed for SOTA parity. $\sim$2 months.
  \item \textbf{Stage 5 (deployment)}: PyO3 bindings; ONNX-equivalent
    serialisation; benchmarks + paper. $\sim$1 month.
\end{enumerate}

Total: $\sim$8 months for a full paradigm-native stack with PyTorch-parity
training and GPU acceleration.

\section*{Rollback path}

Each stage is additive. Falling back to PyTorch is always available
because HSiKAN's existing PyTorch implementation stays the research
path. HymeKo-Nagare is the deployment / inference / paper-systems-track
companion, not a replacement during the transition.

\section*{Risk anticipation}

\begin{itemize}
  \item \textbf{Numerical drift from PyTorch}: Catmull-Rom in pure Rust
    may have different FP rounding than PyTorch's CUDA path; verified
    parity test in every op-pair test.
  \item \textbf{Atomic-add contention}: high-degree hub vertices receive
    many gradient updates from cycles incident to them. Possible
    contention hotspot; mitigated with per-thread accumulators
    combined at end of step.
  \item \textbf{GPU bandwidth-bound on Epinions}: scatter is
    bandwidth-bound on big graphs; wgpu may not match cudarc here.
    Acceptable: cudarc as the fast path; wgpu as portable fallback.
  \item \textbf{Lost research velocity}: Rust's compile-test loop is
    slower than PyTorch eager. Mitigate by keeping research in PyTorch
    (Stage 1's inference-only scope) until Stage 2 lands.
  \item \textbf{Ecosystem cost}: no HF, no torch.compile, no FlashAttention.
    Acceptable for our specific signed-hypergraph model class; revisit
    if scope expands to attention-heavy transformers.
\end{itemize}

\section*{Pitch}

\begin{quote}
\emph{``Signed-hypergraph learning admits a fundamentally different
computational paradigm than dense-tensor DAG autograd: cycle pools are
sums of independent per-cycle terms; gradients decompose the same way.
This factors into MapReduce / dataflow rather than a serial tensor-op
graph, unlocking total parallelism over cycles with no batch-dimension
bottleneck. We implement HymeKo-Nagare as a SoA Rust stack with per-primitive
forward+backward kernels, lockless atomic gradient accumulation, and
Rayon/wgpu execution backends --- yielding $22\times$ over PyTorch eager
on Bitcoin OTC and scaling linearly with core count.''}
\end{quote}

\textbf{Target venue}: MLSys (\textit{The Conference on Machine
Learning and Systems}), 2027 cycle. Companion: ASPLOS or PPoPP for
the systems angle; ICLR if the paradigm-shift framing finds a reviewer
match.

\end{document}
